% Laget automatisk av elevark/lag_elevark.py fra sim/trappelop.html. Endringer her blir overskrevet.
\documentclass[11pt]{article}
\usepackage{elevark}
\begin{document}
\elevarkhode{Fysikk 1 · Mekanisk energi}{Trappeløp}{https://lektorbo.no/sim/trappelop.html}{Simuleringen}

Hvor stor effekt yter du når du løper opp en trapp? Mål høyden på trappa, ta tiden og regn ut arbeidet du gjør mot tyngden. Samle målingene til hele klassen i tabellen nederst.

\section*{Oppgaver}

Dere trenger en trapp, et målebånd og en stoppeklokke. Løp bare så fort som det er trygt.

\begin{oppgaveliste}
\oppgave \fase{gjett}{Gjett}Hvor stor effekt tror du at du yter når du løper opp trappa? Mer eller mindre enn en vannkoker på 2000~W? Mer eller mindre enn én hestekraft?
\svar
\oppgave \fase{test}{Mål}Mål høyden på ett trinn og tell trinnene. Ta tiden når du går opp i vanlig tempo, og når du løper så fort du trygt kan.
\svar
\oppgave \fase{test}{Regn}Regn ut arbeidet du gjør mot tyngden og effekten i begge tilfellene. Sjekk svarene med kalkulatoren. Er arbeidet forskjellig når du går og når du løper? Hva med effekten?
\svar
\oppgave \fase{forklar}{Forklar}Kroppen din bruker mye mer energi enn mgh på å komme opp trappa. Hvor blir resten av energien av?
\svar
\oppgave \fase{test}{Regn}En sjokolade inneholder 2200~kJ. Kroppen omdanner ca. 25~\% av energien fra maten til mekanisk arbeid. Hvor mange ganger må du løpe opp trappa for å bruke opp energien i sjokoladen?
\svar
\oppgave \fase{utfordring}{Utfordring}Sammenlign målingene i klassetabellen. Hvorfor kan det være mer rettferdig å sammenligne effekt per kilo enn effekt? Hvem i klassen ville kommet best ut i et sykkelritt opp en lang bakke?
\svar
\end{oppgaveliste}

\end{document}
